\section{Reception: Hammer, Bogeyman, Doctor}
\label{sec:reception}

\subsection{The controversialist received}

The \work{Disputationes} organized Catholic controversial theology
for two centuries and organized its opposition as well: the 1907
Catholic Encyclopedia's report that Protestant faculties founded
``special chairs\ldots\ to provide replies to it'' is the classic
measure, and the answering literature (its scale may be gauged
from the polemics catalogued under his name in Sommervogel's
Jesuit bibliography, to which the encyclopedia refers readers)
kept his name a fixture of Protestant reading lists long after
his death. Within the Church his authority on the
primacy proved durable in a specific, documentable way: the 1931
doctoral letter states that the Fathers of the Vatican Council
(1870) ``made the greatest use of his writings and opinions''
in the matter of the infallible magisterium (AAS~23, 436; M). His
catechism, continuously reprinted, remained the ordinary food of
Italian catechesis into the twentieth century (AAS~23, 436; M);
his retreat books crossed confessional lines in translation. The
reader who wants the works' afterlife in one sentence may take
the doctoral letter's: he was ``held as a Doctor of the Church
and reverently invoked'' by theologians long before the title was
conferred (AAS~23, 436--437; M).

\subsection{The adversarial memory}

The hostile reception is as old and as continuous, and it has
three distinguishable layers, which this study keeps apart.
First, the confessional and regalist enmity of his own age: the
English and Gallican polemic, the Parlement's censures, the
Venetian pamphlets---an enmity of adversaries who took his
arguments seriously enough to censure them. Second, the
intra-Catholic resistance documented in his cause: the vetoes of
1677 and 1753, argued from his self-memoir and his political
theology, and pressed by reasons of state (\S\ref{sec:cause}).
Third, the modern secular memory, formed in the nineteenth
century around two names: Bruno, executed on the Campo de' Fiori
by sentence of the tribunal in which Bellarmine sat, and Galileo,
whose Vatican dossier, published from the 1870s (Berti,
L'Epinois, von Gebler, then Favaro's national edition), made
Bellarmine's warning and certificate the most examined two pages
of his life. In that literature Bellarmine became a symbol---for
some of obscurantism, for others of a caution more scientific
than Galileo's certainty. Both symbols outrun the documents, as
\S\ref{sec:galileo} sets out; the documents leave a churchman who
enforced a theological consensus later abandoned by the Church
itself, and who stated, conditionally, the principle by which it
would be abandoned.

His relation to religious coercion belongs to this ledger and is
stated within this study's evidence: he served for a quarter
century, as consultor and cardinal, the tribunals whose coercive
character needs no argument; and the received account of his
theology---that the \work{Disputationes} defend the punishment of
obstinate heretics by the secular arm, a common doctrine of his
age---is consistent with everything read here, though the exact
loci (in the controversy on the Church's members) were not
collated for this study and the claim is carried at reported
level. The canonization of 1930 did not canonize any such
doctrine---an act's object is the person's heroic virtue and the
approved miracles, not each opinion---and the honest historical
statement runs in both directions: he cannot be modernized into a
tolerationist, and the commonness of the doctrine in his century
does not remove his documented service of that century's coercive
machinery (S/K, at the stated ceiling).

\subsection{The rehabilitations}

The official twentieth-century acts were themselves a
re-evaluation---the 1930 miracle decree opens by acknowledging
``how many and how great the assaults'' on his memory (AAS~22,
278; M)---and they were accompanied by a scholarly one: the
documentary publications and critical lives of the early
twentieth century, and the standard modern biography by James
Brodrick (named here but not consulted: it remains in
copyright). The
present ecclesial statement of his meaning is Benedict~XVI's 2011
catechesis: a portrait organized not around the controversies but
around the spiritual works, closing on the claim that Bellarmine
teaches ``with the example of his own life that there can be no
true reform of the Church unless there is first our own personal
reform and the conversion of our own heart'' (M, Vatican English).
The distance between that portrait and the seventeenth-century
\latin{haereticorum malleus} is itself the best summary of his
reception: each age has canonized the Bellarmine it needed, and
the documents, patient of neither simplification, remain.
